\changecolours{ScoutGreen}
\section[Laboratório Computacional]{Laboratório Computacional: Validação e Pré-Mapeamento}

% ------------------------------------------------------------------------------
\twocolframe[title={Estudo de Caso: NOC Datacenter IFCE Tauá},leftcol=7cm,rightcol=7cm]{%
  \alert{O Minimundo Operacional}\par
  \itemseps[topsepi=1mm,itemsepi=1.5mm,labelsep=2mm,leftmargini=3.5mm]
  \begin{itemize}
    \item \textbf{Cenário Real:} O campus opera infraestrutura com roteadores core, switches de distribuição e firewalls.
    \item \textbf{Demanda do Sistema:} Registrar cada ativo com identificação patrimonial, IP de gerência, localização precisa no rack e lista de VLANs autorizadas.
    \item \textbf{Atributo Derivado:} Obter o tempo de atividade em dias sem precisar atualizar o banco diariamente.
  \end{itemize}
}{%
  \alert{Mapeamento para o MER}\par
  \itemseps[topsepi=1mm,itemsepi=1.5mm,labelsep=2mm,leftmargini=3.5mm]
  \begin{itemize}
    \item \textbf{Entidade:} \texttt{DISPOSITIVO\_REDE}
    \item \textbf{Identificador:} \underline{\texttt{patrimonio\_id}} (Chave simples)
    \item \textbf{Simples:} \texttt{hostname}, \texttt{ip\_loopback}, \texttt{fabricante}, \texttt{data\_instalacao}
    \item \textbf{Composto:} \texttt{localizacao\_rack} (\texttt{bloco}, \texttt{sala}, \texttt{rack}, \texttt{posicao\_u})
    \item \textbf{Multivalorado:} \texttt{vlan\_tags} ($\{10, 20, 30, \dots\}$)
    \item \textbf{Derivado:} \texttt{dias\_operacao} ($\texttt{hoje} - \texttt{data\_instalacao}$)
  \end{itemize}
}

% ------------------------------------------------------------------------------
\begin{frame}[fragile]{Metamodelo em Python: Validação Conceitual}
  \vspace{0.2cm}
  \begin{columns}[t]
    \begin{column}{0.48\textwidth}
      \footnotesize
      \alert{Definição do Esquema Conceitual}\par
      \vspace{0.15cm}
      \begin{itemize}\setlength{\itemsep}{1.5mm}\setlength{\parskip}{0pt}
        \item Metamodelo orientado a objetos em Python.
        \item Validação de unicidade para a chave identificadora.
        \item Cálculo dinâmico do atributo derivado em tempo de execução.
      \end{itemize}
    \end{column}

    \begin{column}{0.48\textwidth}
      \alert{Código do Metamodelo}\par
      \vspace{0.15cm}
\begin{lstlisting}[language=Python,basicstyle=\ttfamily\fontsize{5.5}{6.5}\selectfont]
esquema = EsquemaEntidade("DISPOSITIVO_REDE")
esquema.adicionar(DefAtributo("patrimonio_id", Tipo.IDENTIFICADOR))
esquema.adicionar(DefAtributo("hostname", Tipo.SIMPLES))
esquema.adicionar(DefAtributo("localizacao_rack", Tipo.COMPOSTO))
esquema.adicionar(DefAtributo("vlan_tags", Tipo.MULTIVALORADO))
esquema.adicionar(DefAtributo("dias_operacao", Tipo.DERIVADO,
                             funcao=lambda d: (date.today() - d["data_inst"]).days))
\end{lstlisting}
    \end{column}
  \end{columns}
\end{frame}

% ------------------------------------------------------------------------------
\begin{frame}[fragile]{Ponte para o Modelo Relacional: O Impacto da 1FN}
  \vspace{0.2cm}
  \begin{columns}[t]
    \begin{column}{0.48\textwidth}
      \footnotesize
      \alert{A Regra da 1ª Forma Normal (1FN)}\par
      \vspace{0.15cm}
      \begin{itemize}\setlength{\itemsep}{1.5mm}\setlength{\parskip}{0pt}
        \item Proíbe atributos compostos e multivalorados em colunas únicas!
        \item \textbf{Composto:} Planificado em colunas atômicas (\texttt{rack\_bloco}, \texttt{posicao\_u}).
        \item \textbf{Multivalorado:} Exige \textbf{tabela auxiliar}.
        \item \textbf{Derivado:} Calculado dinamicamente via \textbf{VIEW}.
      \end{itemize}
    \end{column}

    \begin{column}{0.48\textwidth}
      \alert{Esquema Relacional em SQLite}\par
      \vspace{0.15cm}
\begin{lstlisting}[language=SQL,basicstyle=\ttfamily\fontsize{5.5}{6.5}\selectfont]
CREATE TABLE dispositivos_rede (
    patrimonio_id TEXT PRIMARY KEY,
    hostname TEXT NOT NULL UNIQUE,
    data_instalacao DATE NOT NULL,
    rack_bloco TEXT NOT NULL,
    rack_posicao_u TEXT NOT NULL
);

CREATE TABLE dispositivo_vlans (
    id INTEGER PRIMARY KEY,
    dispositivo_id TEXT NOT NULL,
    vlan_id INTEGER NOT NULL,
    FOREIGN KEY(dispositivo_id)
      REFERENCES dispositivos_rede(patrimonio_id)
);
\end{lstlisting}
    \end{column}
  \end{columns}
\end{frame}
